\subsection{Design}
Although \vts could be identified relatively well by learning the rules based on the frequency of code tokens, not all the parts of a vulnerable function are important for determining the type of vulnerabilities in it. 
%
%
There are certain code elements which are more crucial than the others in determining the types of a vulnerability. For example, \textit{Buffer Overflow} vulnerabilities usually cause by the \textit{assignments} (e.g., assigning too large index) or \textit{function calls} (e.g., copying data larger a buffer's capacity). In Big-Vul dataset, about 95\% of vulnerable statements are/contain either \textit{function calls} (78\%), \textit{assignments} (44\%), \textit{control structures} (38\%), or \textit{return statements} (19\%). Instead of considering all code elements equally important, focusing on those crucial kinds of code elements in the vulnerable functions could improve the accuracy of VTI predictions.
%
% \textbf{Attention.} Criticize other attentions. Extract distinguishing token in a selective locations. 
To realize the above idea, attention mechanism~\cite{attention} could be a great option. However, this mechanism has its own drawbacks. This method adds more weight parameters to the neural network model. This could significantly increase training time. Furthermore, attention mechanism are applicable for NN models, but not traditional models. 


To cope with these problems, we design a light-weight technique which can be independently applied as a post-processing step for every VTI approach. 
%
Our idea is that the code tokens, which can be used to distinguish each \vt from the others (so called, \textit{distinguish tokens}), are identified by statistically analyzing the code elements critical for VTI such as \textit{function calls}, \textit{assignments}, \textit{control structures}, and \textit{return statements}.  
%
These distinguishing tokens are used to adjust the predictions produced by VTI approaches. Particularly, a function $f$ should or should not be of a type $t$ if $f$ has the distinguishing tokens of $t$.


\begin{definition}{\textit{\textbf{(Distinguishing token).}}}
For a type $t \in T$, a distinguishing token in a code element is a code token which are more/less \textit{prevalent} in the code element of the functions having vulnerabilities of type $t$ than in any others.
\end{definition}
%
The prevalence of a token $c$ regarding a code element $e$ in the cases of type $t$, $D_t$, ($prev(c, e, D_t)$), is reflected via the ratio of the code elements $e$ of cases containing $c$ in $D_t$. Formally $prev(c, e, D_t) = \frac{count(c, e, D_t)}{|D_t|}$, where $count(c, e, D_t)$ is the number of cases where the elements $e$ contain $c$. Token $c$ is a positive distinguishing tokens of $t$, $c \in \Delta^+(t, e)$, if $c$ is more prevalent in the cases of type $t$ than any other types: 
%
$$
dis^+(c, e, t, D) = \frac{prev(c, e, D_t)}{\max_{t' \in T \setminus \{t\}} prev(c, e, D_{t'})} > 1
$$
%
%
When $\max_{t' \in T \setminus \{t\}} prev(c, e, D_{t'}) = 0$, then $dis^+(c, e, t, D)$ is infinity. In that case, $c$ might only appear in $D_t$, not in the others.
%
%

Meanwhile, regarding a code element $e$, there are certain tokens which are more prevalent in the cases of the other types rather than $t$ . Such code tokens are considered as negative distinguishing tokens of $t$ regarding $e$, $c \in \Delta^-(t, e)$. Formally, $c \in \Delta^-(t, e)$ if $c$ satisfies the following condition:
% when $c$ is less prevalent in $D_t$ than in the cases where $c$ is least prevalent:
%
% $$
% dis^-(c, e, t, D) = \frac{prev(c, e, D_t)}{\min_{t' \in T \setminus \{t\}} prev(c, e, D_{t'})} < 1
% $$
$$
dis^-(c, e, t, D) = \frac{\min_{t' \in T \setminus \{t\}} prev(c, e, D_{t'})}{prev(c, e, D_t)} > 1
$$
%
%
%
% Figure X shows the numbers of positive and negative distinguishing tokens of the popular types in a training dataset (8/10 random folds of Big-Vul). 
%
If $c$ is never in any case of type $t$, then $prev(c, e, D_t) = 0$ and $dis^-(c, e, t, D) = \infty$.
%
Intuitively, regarding a code element, \textit{when a case has a positive distinguishing token of type $t$, then the vulnerable function are likely to have a vulnerability of $t$} (\rp). Similarly, \textit{if the case has a negative distinguishing token of type $t$, the vulnerable function are likely to be a case of any types other than $t$} (\rn).


A vulnerable function $f$ might satisfy neither \rp nor \rn. The prediction for $t$ has to rely on the prediction of a VTI approach. Thus, we apply \rp and \rn as to design a technique which can be used as \textit{\textbf{an independent step}} combined with the VTI model. Particularly, for the prediction $Z$ of $f$ produced by the model, we apply \rp and \rn to adjust $Z$ to produce the final prediction $Z'$. Particularly, for a type $t \in T$, $Z'[t]$, which is the element corresponding $t$ in $Z'$, is adjusted based on that of $Z$ (i.e., $Z[t]$):
\begin{itemize}
    \item If $Z[t] = 0$, yet a code element $e$ of $f$ contains $c \in \Delta^+(t, e)$, then $Z'[t] = 1$.
    \item If $Z[t] = 1$, yet a code element $e'$ of $f$ contains $c' \in \Delta^-(t, e)$, then $Z'[t] = 0$.
    \item Otherwise, $Z'[t] = Z[t]$.
\end{itemize}